Use the sky map Map-OM1 to answer the questions below. The map has a stereographic projection, and planets and the Moon are present on the map.
\begin{description}
\item[(OM1.1)] Mark all the stars listed below with a square ($\square$) drawn around each star on Map-OM1. Label every marked star with its corresponding object number. \hfill[10]
\begin{center}
\begin{tabular}{rllrll}
1 & $\alpha$ Ari & Hamal & 6 & $\beta$ CMi & Gomeisa\\
2 & $\delta$ Cyg & Fawaris & 7 & $\alpha$ Lep & Arneb\\
3 & $\beta$ Dra & Rastaban & 8 & $\delta$ UMa & Megrez\\
4 & $\alpha$ Cet & Menkar & 9 & $\beta$ Tau & Elnath\\
5 & $\delta$ Eri & Rana & 10 & $\alpha$ And & Alpheratz\\
\end{tabular}
\end{center}
\item[(OM1.2)] Use the provided plus signs ($+$) on the map and label them corresponding to the following eight Messier objects, M15, M31, M34, M35, M39, M42, M44, and M45. \hfill[8]
\item[(OM1.3)] The map is centered at the zenith. Draw the meridian and then draw with star-line patterns the constellations along or near the meridian and note the IAU abbreviations on the sheet. \hfill[10]
\item[(OM1.4)] Draw the celestial equator on Map-OM1. Mark it as ``EQ''. \hfill[2]
\item[(OM1.5)] Mark the cardinal points as N, S, W, and E on Map-OM1. \hfill[2]
\item[(OM1.6)] Draw the ecliptic on Map-OM1 and label it ``EC''. \hfill[2]
\item[(OM1.7)] Draw the galactic equator on Map-OM1 and label it ``GE''. Mark either the galactic center (GC) or anticenter (GA), whichever is visible on the map. \hfill[4]
\item[(OM1.8)] Draw the circumpolar circle on Map-OM1 and label it ``CC''. \hfill[2]
\end{description}
\ioaafig{2025_OBS-SKY_Q10-f1.pdf}
